\section[Weighted corridor routing]{Charged corridor routing for an arbitrary supplied tree}
\label{sec:op3-weighted-corridor-routing}

\paragraph{Input and scope.}
The graph in this section is supplied in full, connected, loopless and
positively weighted; parallel edges are permitted as separate records.
A spanning tree $T$ is also supplied by original-edge indices. No
low-stretch property or weight-ratio bound is assumed. Write
\[
 r_e=1/c_e,\qquad w_e=c_e r(T(e)),\qquad W=\sum_{e\in E(G)}w_e.
\]
For $n\geq2$, $W\geq n-1$, because every retained tree edge contributes
one. The following construction supplies the missing routing step for
\cref{thm:op3-rooted-piece-spectral-comparison}; selection of $T$ and
fast core sparsification remain separate source obligations.

\begin{lemma}[Corridor cuts, roots and the common offset; \statusproved]
\label{lem:op3-corridor-common-offset}
Suppose connected pieces partition $E(T)$ and cover $V$, distinct pieces
meet in at most one vertex, and every piece has at most two shared
boundary vertices. Every vertex is owned by one containing piece.
Let $C$ be all shared vertices, or one arbitrary vertex if there is a
single piece. In every two-boundary piece $W_i$, delete an edge $f_i$
minimizing the global path congestion
\[
 g(f)=\sum_{e:f\in T(e)}c_e
\]
on the boundary-to-boundary corridor $R_i=T(s_i,t_i)$.
Keep every other tree edge. The resulting forest has exactly one
vertex of $C$ per component. The connected parts within individual
pieces are rooted pieces in the sense of
\cref{sec:op3-rooted-piece-comparison}, allowing empty-edge root parts.

Let $D_i$ be the original edges assigned to $W_i$ by endpoint ownership
whose tree path meets an edge of $W_i$. On every $f\in R_i$,
\[
 g(f)=g_i(f)+b_i,\qquad
 g_i(f)=\sum_{e\in D_i:f\in T(e)}c_e,
\]
where $b_i\geq0$ is independent of $f$. Consequently $f_i$ also
minimizes $g_i$ on $R_i$.
\end{lemma}
\begin{proof}
With a single piece, make no cut and use the chosen root; the forest
assertion is immediate. There is no two-boundary corridor to analyze.
If there is more than one piece, every piece has a shared boundary.
Cutting a two-boundary piece on its corridor produces two connected
parts, one containing each boundary; a one-boundary piece is already
rooted. Distinct original pieces meet only at shared vertices.
All parts therefore meet only at their common root. Every vertex is
covered by such a part, and parts with distinct roots cannot connect.
This proves the forest and rooted-piece assertions, including singleton
owner pieces that contain no tree edge.

An interior vertex of $W_i$ belongs to no other piece, so its owner is
$W_i$. Also, an edge of $T$ leaving $W_i$ can do so only at a shared
boundary, since its other containing piece then meets $W_i$ there.
Thus an original edge not assigned to $W_i$ has no interior endpoint.
If its tree path crosses any edge of $R_i$, it must enter and leave
$W_i$ at the two distinct boundaries, possibly starting or ending at
one of them. It traverses the entire corridor. Edges assigned to
$W_i$ but outside $D_i$ cross no edge of $W_i$ and contribute zero.
The sum of the conductances of the remaining unassigned crossing edges
is the common offset $b_i$. Adding that constant preserves all
minimizers, including deterministic edge-index tie breaking.
\end{proof}

\begin{lemma}[Routed stretch from owned loads; \statusproved]
\label{lem:op3-corridor-owned-stretch}
Under \cref{lem:op3-corridor-common-offset}, every rooted part $B$
of a piece $W_i$ satisfies
\[
 \sum_e c_e r(P_F(e)\cap B)
 \leq 2\sum_{e\text{ assigned to }W_i}w_e.
\]
Here $P_F(e)$ is the original-edge route defined in
\cref{sec:op3-rooted-piece-comparison}. The assertion includes original
edges retained in $F$ and pieces containing many branches at one root.
\end{lemma}
\begin{proof}
An edge contributes a forest segment within $W_i$ only if at least one
of its endpoints is a nonroot interior vertex of $W_i$. A route from
another piece terminates when it reaches the shared root and cannot
continue through $W_i$. Such a contributing original edge is assigned
to $W_i$, and its tree path meets $W_i$, so it belongs to $D_i$.

For a one-boundary piece, the forest segment is contained in $T(e)$,
which already gives the stronger factor-one bound. Consider a
two-boundary piece and write $J=T(e)\cap W_i$ as an edge path.
If $f_i\notin T(e)$, routing adds no edge within $W_i$: the segment
is contained in $J$. If $f_i\in T(e)$, the endpoint segments in the
two cut sides can instead run toward the two boundary roots. Their
edges are contained in $J\cup(R_i\setminus J)$. To see this, project
each endpoint of $J$ onto the corridor. The paths off the corridor
remain unchanged, and the two paths toward the boundary roots use
only corridor edges outside the interval between those projections.
This also covers a projected endpoint already equal to a boundary.
It follows in both cases that
\[
 r(P_F(e)\cap W_i)
 \leq r(T(e))+
       \boldsymbol{1}_{\{f_i\in T(e)\}}r(R_i\setminus T(e)).
\]
The left side sums the forest portions in both rooted parts of $W_i$;
no core edge is included. The additional term is present only when
the deleted edge lies on the original tree path, correcting the
reversed displayed cases in the source's Claim 5.14.

Discarding the nonpositive subtraction from the extra length gives
\[
 \sum_e c_e r(P_F(e)\cap W_i)
 \leq\sum_{e\in D_i}w_e+r(R_i)g_i(f_i).
\]
By the preceding lemma, $f_i$ minimizes $g_i$, and hence
\[
 \begin{split}
 r(R_i)g_i(f_i)
 &\leq\sum_{f\in R_i}r_f g_i(f)\\
 &=\sum_{e\in D_i}c_e r(T(e)\cap R_i)
 \leq\sum_{e\in D_i}w_e.
 \end{split}
\]
All terms are nonnegative. Restricting to either rooted part and
enlarging $D_i$ to all assigned edges proves the claim. A singleton
piece has no forest segment and satisfies it trivially.
\end{proof}

\begin{theorem}[Weighted construction from a supplied spanning tree; \statusproved]
\label{thm:op3-supplied-tree-constructor}
For a supplied connected positive-weight graph on $n\geq2$ vertices,
a supplied spanning tree of total stretch $W$, and $1\leq j\leq n$,
one can construct a rooted subforest and an aggregated core on at most
$j$ roots such that their weighted union $Q$ satisfies
\[
 L(G)\preceq L(Q)\preceq \frac{5376W}{j}L(G).
\]
The construction costs $O((n+m)\log(2+n+m))$ exact-real word work and
$O(n\log(2+n)+m)$ total allocated words, including all input reads,
tree indexing, decomposition, root labeling, sorting and output.
The core can have $O(m)$ edges; it has not yet been sparsified.
There is no minimum-weight or weight-ratio assumption.
\end{theorem}
\begin{proof}
For $j<64$, use all of $T$ as one rooted piece. Its total local stretch
is $W$. Otherwise set $\theta=64W/j$, assemble
$a_v=\sum_{e\ni v}w_e$, and use
\cref{thm:op3-weighted-tree-ownership,cor:op3-weighted-edge-ownership}.
There are at most $j/4$ pieces, with assigned stretch below
$2\theta$ for each nontrivial piece. Use the corridor cuts above.
If there are $h>1$ pieces, their number of shared vertices is at most
$h-1$ by the incidence count in the ownership proof. If $h=1$, use
one root. Thus there are at most $j$ roots in all cases.

Every nonempty rooted piece has local stretch below $4\theta$ in
the large-$j$ branch. Set $\kappa=256W/j$. Since $W\geq n-1$,
$\kappa\geq1$, and it also exceeds $W$ in the small-$j$ branch.
\Cref{thm:op3-rooted-piece-spectral-comparison} constructs $Q$ by
scaling forest edges by $3\kappa$ and aggregated core conductances
by three. Its comparison gives $21\kappa=5376W/j$ as claimed.

For completeness, tree paths are not supplied for free. Build parent,
depth and resistance-prefix arrays and a binary-ancestor table in
$O(n\log(2+n))$ work and allocated words. Each original edge uses
one lowest-common-ancestor query to compute its tree resistance.
For its conductance $c_e$, add $c_e$ at both endpoints and subtract
$2c_e$ at that ancestor. A reverse tree traversal sums these differences;
the sum below every parent edge is exactly its global congestion.
All original stretches, vertex loads and congestion values are therefore
computed in $O((n+m)\log(2+n))$ work without a complete path table.

The ownership and refinement work is linear in $n$. Enumerate a
two-boundary corridor by ascending its two parent chains to their
lowest common ancestor. Corridors in distinct pieces have disjoint
tree edges, so total path output and minimum-congestion scanning cost
$O(n)$, besides $O(h\log(2+n))$ ancestor queries. Rooted-part adjacency
is assembled from each piece's own edge list, never by repeatedly
reading a shared root's complete adjacency. These lists, all local
traversals, and the global forest root labels have total linear size.

Finally scan the original edges, sort cross-root records by their
two root labels, add parallel conductances, and emit the scaled forest
and core. A merge sort uses one reusable $O(m)$ buffer. Ancestor
tables, initial and refined piece lists, linked-list arenas, local
adjacencies, released scratch arrays and emitted records all fit the
stated total-allocation bound. Exact arithmetic is charged as word
work; this is not a bit-complexity assertion.
\end{proof}

For $n=1$, the graph has no edges and the empty graph is an exact
comparison with one root. This case is handled directly and never
divides by $W=0$.

\paragraph{Audit.}
The implementation constructs all of the above objects from original graph records and
supplied tree indices. Independent tree paths, congestion sums, rooted
part intersections, corrected routing inequalities and dense PSD tests
serve only as validators. The final exact audit and its source/backend
hashes are recorded in \texttt{WEIGHTED\_CORRIDOR\_ROUTING\_AUDIT.json}.
The source is \texttt{weighted\_corridor\_routing.py} in the note's
proof-audit directory.
The full audit passes 27,961 supplied constructions, 55,706 exact PSD
certificates, 33,404 common-offset corridor checks and 16,189 cases
in which the deleted edge lies on an original path. It checks every
spanning tree and every root of the connected atlas graphs through
five vertices, two weight profiles, and five load thresholds, as well
as the actual $j$-budget branch. Larger stars, paths and random trees
through 512 vertices include parallel original edges and weights at
$2^{-80}$ and $2^{80}$. Independent validators check 272,880 stretches,
134,700 congestion values and 132,199 rooted-piece loads. These finite
checks support the stated proof and do not replace its asymptotic work
argument.
This theorem does not provide a low-stretch tree, a sparse core, or
local discovery. Reconciling the arbitrary-weight source primitives and
their requested confidence remains necessary for the full supplied
near-linear recurrence. General OP3 remains \statusopen.
